% Unit 087 terjemahan Bahasa Indonesia dari chapter6.tex baris 2928--3141.
% Sumber: Methods of Algebra, Volume 2, commit
% 9a5803ff77dd3257484cb177f851a73770a59dd3 (CC BY 4.0).
% stable-unit-id: o014.aljabr2.chapter6.exercises
% Cakupan: seluruh latihan Bab 6; menutup Bab 6.

% segment-id: o014.li2.chapter6.exercises.h001
\begin{Exercises}
	% segment-id: o014.li2.chapter6.exercises.ex001
	\item Misalkan $\Bbbk$ gelanggang komutatif dan $t\in\Bbbk$. Untuk setiap
	modul-$\Bbbk$, tuliskan $m_t$ untuk endomorfisme perkalian dengan $t$.
	Misalkan $M$ modul-$G$. Tunjukkan bahwa endomorfisme yang diinduksi oleh
	$m_t\in\End_G(M)$ pada $\Hm^n(G,M)$ dan $\Hm_n(G,M)$ juga tepat merupakan
	$m_t$.

	\begin{hint}
		Untuk kohomologi, tuliskan endomorfisme terinduksi pada kompleks standar,
		atau ambil resolusi injektif $0\to M\to I^0\to\cdots$ dan tunjukkan
		bahwa endomorfisme terinduksi berasal dari diagram
		\[\begin{tikzcd}
			0 \arrow[r] & M \arrow[r] \arrow[d,"{m_t}"] &
			I^0 \arrow[r] \arrow[d,"{m_t}"] &
			I^1 \arrow[r] \arrow[d,"{m_t}"] & \cdots\\
			0 \arrow[r] & M \arrow[r] & I^0 \arrow[r] & I^1 \arrow[r] & \cdots.
		\end{tikzcd}\]
	\end{hint}

	% segment-id: o014.li2.chapter6.exercises.ex002
	\item Misalkan modul-$G$ $M$ projektif sebagai modul-$\Bbbk$. Buktikan
	bahwa $\mathsf{L}\otimes M\to\Bbbk\otimes M\simeq M$ memberikan resolusi
	projektif modul-$G$ $M$; pernyataan yang sama berlaku jika $\mathsf{L}$
	diganti dengan $\overline{\mathsf{L}}$.

	\begin{hint}
		Terapkan Proposisi~\ref{prop:group-tensor-proj}.
	\end{hint}

	% segment-id: o014.li2.chapter6.exercises.ex003
	\item Buktikan Proposisi~\ref{prop:group-UCT} memakai teorema koefisien
	universal untuk kohomologi kompleks rantai.

	% Another candidate: Rotman, Homological Algebra. Duality Theorem 9.73. Think about it.

	% segment-id: o014.li2.chapter6.exercises.ex004
	\item Misalkan $p$ bilangan prima, $G$ grup-$p$, dan $M$ modul-$G$ yang
	memenuhi $pM=0$. Buktikan bahwa sifat-sifat berikut ekuivalen:
	\begin{inparaenum}[(i)]
		\item $M=0$,
		\item $M^G=0$,
		\item $M_G=0$.
	\end{inparaenum}

	\begin{hint}
		Untuk membuktikan (ii) $\implies$ (i), ambil
		$x\in M\smallsetminus\{0\}$ dan perhatikan subruang vektor-$\F_p$ $E$
		dari $M$ yang dibangun oleh $\{gx:g\in G\}$. Gunakan
		$|E^G|\equiv|E|\pmod p$ untuk menunjukkan $M^G\neq0$.

		Untuk membuktikan (iii) $\implies$ (i), ambil
		$M^\vee:=\Hom_{\F_p}(M,\F_p)$ dengan struktur modul-$G$ alami.
		Perhatikan
		$(M^\vee)^G\simeq\Hom_{\F_p}(M_G,\F_p)$, lalu deduksikan
		$M_G=0\implies M^\vee=0\implies M=0$.
	\end{hint}

	% segment-id: o014.li2.chapter6.exercises.ex005
	\item Misalkan $F$ suatu medan. Sebuah \emph{representasi} grup $G$ adalah data
	$(\rho,V)$, dengan $V$ ruang vektor-$F$, $\GL(V)$ grup automorfisme
	linearnya, dan $\rho:G\to\GL(V)$ homomorfisme grup. Jelas ini sama dengan
	modul-$G$ berkoefisien $F$. Sebuah \emph{representasi projektif} adalah
	data $(\pi,V)$ dengan homomorfisme
	$\pi:G\to\PGL(V):=\GL(V)/F^\times\identity$. Isomorfisme antarrepresentasi
	representasi atau representasi projektif didefinisikan dengan cara yang
	jelas. Selanjutnya identifikasikan $F^\times$ dengan
	$F^\times\identity$.
	\index{representasi (representation)}
	\index{representasi projektif (projective representation)}

	\begin{enumerate}[(i)]
		\item Untuk suatu representasi projektif $(\pi,V)$ dari $G$, andaikan
		untuk setiap $g\in G$ dipilih prapeta
		$\widetilde\pi(g)\in\GL(V)$ dari $\pi(g)$. Periksa bahwa
		\[ \widetilde\pi(g_1)\widetilde\pi(g_2)
		=c(g_1,g_2)\widetilde\pi(g_1g_2) \]
		menentukan $c\in Z^2(G,F^\times)$, dengan aksi $G$ pada $F^\times$
		trivial. Jika $\widetilde\pi(1_G)=\identity_V$, maka
		$c\in\overline{Z}^2(G,F^\times)$.
		\item Buktikan bahwa citra $c$ dalam $\Hm^2(G,F^\times)$ hanya bergantung
		pada kelas isomorfisme $(\pi,V)$ dan tidak bergantung pada pilihan apa
		pun.
		\item Perhatikan ekstensi pusat grup
		$1\to F^\times\to\widetilde G\to G\to1$
		(Definisi~\ref{def:central-ext}), dengan
		\[ \widetilde G:=\{(g,\widetilde\pi)\in G\times\GL(V):
		\pi(g)=\widetilde\pi\bmod F^\times\}. \]
		Homomorfisme pertama memetakan $t\in F^\times$ ke $(1_G,t)$ dan yang
		kedua adalah proyeksi. Tunjukkan bahwa ekstensi ini isomorfik dengan
		ekstensi pusat yang ditentukan oleh kelas kohomologi dalam~(ii).
		\item Berikan bijeksi
		\[ \{\text{representasi }(\sigma,V)\text{ yang mengangkat }(\pi,V)\}
		\xleftrightarrow{1:1}
		\{\text{pembelahan ekstensi pusat }\widetilde G\}. \]
		Tunjukkan pula bahwa $(\pi,V)$ dapat diangkat menjadi representasi jika
		dan hanya jika kelas kohomologi dalam~(ii) trivial.
		\item Tuliskan secara eksplisit grup $G$ dan representasi projektif
		$(\pi,V)$ yang tidak dapat diangkat menjadi representasi. Cobalah mencari
		contoh dengan $\pi:G\to\PGL(V)$ taktrivial.
	\end{enumerate}

	% segment-id: o014.li2.chapter6.exercises.ex006
	\item Untuk subgrup $G_1\subset G_2\subset G_3$, buktikan isomorfisme
	kanonik funktor
	$\Ind^{G_3}_{G_2}\Ind^{G_2}_{G_1}\simeq\Ind^{G_3}_{G_1}$ dan
	$\iInd^{G_3}_{G_2}\iInd^{G_2}_{G_1}\simeq\iInd^{G_3}_{G_1}$.
	Deskripsikan isomorfisme tersebut sejelas mungkin.

	% segment-id: o014.li2.chapter6.exercises.ex007
	\item Buktikan bahwa jika $H$ subgrup dari $G$, maka
	$\mathrm{cd}(H)\leq\mathrm{cd}(G)$ dan
	$\mathrm{hd}(H)\leq\mathrm{hd}(G)$; lihat Definisi~\ref{def:group-dim}.
	\begin{hint}
		Terapkan Teorema~\ref{prop:Shapiro}.
	\end{hint}

	% segment-id: o014.li2.chapter6.exercises.ex008
	\item Tunjukkan bahwa untuk $m\in\Z_{\geq1}$ berlaku
	$\mathrm{cd}(C_m)=\mathrm{hd}(C_m)=\infty$. Buktikan bahwa
	$\mathrm{cd}(G)<\infty$ mengakibatkan $G$ bebas torsi, dan bahwa
	$\mathrm{cd}(G)=0$ jika dan hanya jika $G$ grup trivial.
	\begin{hint}
		Untuk bagian pertama, terapkan
		Proposisi~\ref{prop:group-cohomology-cyclic} dengan $A=\Z/m\Z$.
		Untuk bagian lainnya, perhatikan subgrup siklik dari $G$.
	\end{hint}

	% segment-id: o014.li2.chapter6.exercises.ex009
	\item Cobalah membuktikan Korolari~\ref{prop:group-coh-torsion} langsung
	melalui operasi pada kompleks standar.

	% Reference: Cohomology of Number Fields, (1.5.6) and (1.5.7)
	% segment-id: o014.li2.chapter6.exercises.ex010
	\item Misalkan $U,V$ subgrup dari $G$ dan $(G:V)$ hingga. Tuliskan pemetaan
	restriksi kohomologi sebagai $\mathrm{res}^G_U$ dan korestriksi sebagai
	$\mathrm{cor}^V_G$, dan seterusnya; derajat kohomologi $n$ dihilangkan dari
	notasi. Pilih dekomposisi koset ganda
	\[ G=\bigsqcup_{\sigma\in S}U\sigma V, \]
	dengan himpunan hingga $S\subset G$ berupa wakil-wakil koset ganda. Misalkan
	$M$ modul-$G$.
	\begin{enumerate}[(i)]
		\item Untuk setiap $\sigma\in G$ dan $n\in\Z_{\geq0}$, definisikan
		secara wajar
		\[ \mathcal{T}_\sigma:
		\Hm^n(V\cap\sigma^{-1}U\sigma,M)
		\to\Hm^n(U\cap\sigma V\sigma^{-1},M) \]
		sedemikian sehingga untuk $n=0$ pemetaan ini menjadi
		\[ M^{V\cap\sigma^{-1}U\sigma}\rightiso
		M^{U\cap\sigma V\sigma^{-1}},\qquad x\mapsto\sigma x. \]
		\item Buktikan
		\[ \mathrm{res}^G_U\circ\mathrm{cor}^V_G
		=\sum_{\sigma\in S}
		\mathrm{cor}^{U\cap\sigma V\sigma^{-1}}_U
		\circ\mathcal{T}_\sigma
		\circ\mathrm{res}^V_{V\cap\sigma^{-1}U\sigma}. \]
		\item Deskripsikan aksi
		$\mathrm{res}^G_U\circ\mathrm{cor}^U_G$ ketika $U=V\lhd G$.
	\end{enumerate}

	% segment-id: o014.li2.chapter6.exercises.ex011
	\item Misalkan $G$ grup abelian, dipandang sebagai modul-$\Z$. Definisikan
	$G\wedge G$ sebagai hasil bagi $G\dotimes{\Z}G$ oleh submodul yang dibangun
	oleh $\{g\otimes g:g\in G\}$, dan tuliskan $x\wedge y$ untuk citra
	$x\otimes y$. Buktikan $\Hm_2(G,\Z)\simeq G\wedge G$.

	\begin{hint}
		\index[sym1]{[,]}
		Setelah presentasi grup $G=F/R$ ditetapkan, rumus Hopf
		(Korolari~\ref{prop:Hopf-H2}) memberikan
		$\Hm_2(G,\Z)\simeq[F,F]/[F,R]$. Karena $G$ abelian,
		$[F,F]\subset R$. Dari identitas
		$[xy,z]=x[y,z]x^{-1}\cdot[x,z]$, pemetaan
		\[ \psi_0:G\times G\to[F,F]/[F,R],\qquad
		(f_1R,f_2R)\mapsto[f_1,f_2]\bmod[F,R] \]
		bersifat bilinear-$\Z$ dan menginduksi homomorfisme surjektif
		$\psi:G\wedge G\to[F,F]/[F,R]$. Ambil basis $\{x_i\}_{i\in I}$ dari
		$F$ dan beri $I$ urutan total. Kita memerlukan fakta teori grup berikut:
		$H:=[F,F]/[[F,F],F]$ adalah grup abelian bebas dengan basis kelas-koset
		$[x_i,x_j]$ untuk $i<j$. Maka definisikan
		$\phi_0:H\to G\wedge G$ dengan memetakan kelas-koset $[x_i,x_j]$ ke
		$x_iR\wedge x_jR$. Tunjukkan bahwa
		\[ f_1,f_2\in F\implies
		\phi_0(\text{kelas-koset }[f_1,f_2])=f_1R\wedge f_2R. \]
		Dari sini, turunkan $\phi_0$ menjadi homomorfisme
		$\phi:[F,F]/[F,R]\to G\wedge G$, lalu tunjukkan bahwa $\phi$ dan
		$\psi$ saling invers.

		Fakta teori grup di atas untuk $I$ hingga terdapat dalam
		\cite[Teorema 11.2.4]{Ha76}; kasus umum mengikutinya.
	\end{hint}

	% segment-id: o014.li2.chapter6.exercises.ex012
	\item Misalkan $H\lhd G$ dan $M$ modul-$G$. Beri kompleks standar
	$C(H,M^H)$ aksi alami $G/H$ yang menginduksi aksi kanonik $G/H$ pada setiap
	$\Hm^n(H,M^H)$ (lihat uraian sebelum Lema~\ref{prop:LHS-action-aux}).

	% segment-id: o014.li2.chapter6.exercises.ex013
	\item Untuk filtrasi dalam Catatan~\ref{rem:LHS-SS-mult}, deskripsikan
	halaman $E_2$ dari barisan spektralnya dan tunjukkan bahwa sifat-sifat dalam
	Teorema~\ref{prop:LHS-SS} tetap berlaku. Jelaskan kompatibilitas barisan
	spektral dengan hasil kali cawan sejauh mungkin.

	% segment-id: o014.li2.chapter6.exercises.ex014
	\item Perhatikan grup siklik hingga $C_m$ berorde $m$. Untuk setiap
	modul-$C_m$ $A$ dan $n\in\Z_{\geq1}$, buktikan bahwa isomorfisme periodik
	$\Hm^n(C_m,A)\to\Hm^{n+2}(C_m,A)$ dalam
	Proposisi~\ref{prop:group-cohomology-cyclic} sama dengan komposisi dua
	homomorfisme penghubung dari barisan eksak pendek kanonik dalam
	Contoh~\ref{eg:cup-periodicity}:
	\begin{gather*}
		0\to\mathfrak{I}\otimes A\to\Z[C_m]\otimes A\to A\to0,\\
		0\to A\xrightarrow{\nu\otimes\identity}\Z[C_m]\otimes A
		\xrightarrow{(\sigma-1)\otimes\identity}\mathfrak{I}\otimes A\to0.
	\end{gather*}

	\begin{hint}
		Hitung kohomologi dengan kompleks periodik
		$A\xrightarrow{\sigma-1}A\xrightarrow{\nu}A\to\cdots$ dari bukti
		Proposisi~\ref{prop:group-cohomology-cyclic}, dimulai pada suku derajat
		nol, dan gunakan kompleks ini untuk mendeskripsikan homomorfisme
		penghubung. Pisahkan kasus ganjil dan genap; diperlukan beberapa
		perhitungan.
	\end{hint}

	% segment-id: o014.li2.chapter6.exercises.ex015
	\item Misalkan $\Bbbk$ gelanggang komutatif sembarang dan
	$m\in\Z_{\geq1}$. Tentukan struktur aljabar kohomologi
	$\Hm^\bullet(C_m,\Bbbk)$ sebagai aljabar-$\Bbbk$ bergradasi.
%	\begin{hint}
%		Substitusi $A=\Bbbk$ dalam Contoh~\ref{eg:cup-periodicity} mungkin membantu.
%	\end{hint}

	% Reference: Spanier, p.254
	% According to May (p.155), it is rather $\lrangle{\alpha \cup \beta, \gamma} = \lrangle{\beta, \alpha \cap \gamma}$
	% segment-id: o014.li2.chapter6.exercises.ex016
	\item Untuk hasil kali tudung dari Catatan~\ref{rem:cap-product}, buktikan
	relasi adjungsi dengan hasil kali cawan
	\[ \lrangle{\alpha\cup\beta,\gamma}
	=\lrangle{\alpha,\beta\cap x}\in(M_1\otimes M_2\otimes M_3)_G, \]
	dengan $M_1,M_2,M_3$ modul-$G$,
	$\alpha\in\Hm^p(G,M_1)$, $\beta\in\Hm^q(G,M_2)$, dan
	$x\in\Hm_{p+q}(G,M_3)$.

	% segment-id: o014.li2.chapter6.exercises.ex017
	\item Dalam soal ini $G$ grup hingga, gelanggang koefisien dalam teori
	kohomologi dan homologi adalah $\Z$, dan modul-$G$ bebas berarti
	modul-$\Z[G]$ bebas. Tuliskan modul kontragradien dari modul-$G$ $M$ sebagai
	$M^*:=\Hom_{\Z}(M,\Z)$.
	\begin{enumerate}[(i)]
		\item Ambil resolusi bebas $\epsilon:P\to\Z$ dari modul-$G$ trivial,
		dengan $P:=[\cdots\to P_n\to P_{n-1}\to\cdots]$ dan setiap $P_n$
		dibangkitkan secara hingga. Definisikan
		$P^*:=[\cdots\to P_n^*\to P_{n+1}^*\to\cdots]$. Tunjukkan bahwa
		$\epsilon^*:\Z\to P^*$ merupakan kuasi-isomorfisme.
		\begin{hint}
			Setiap $P_n$ adalah modul-$\Z$ bebas. Jika kompleks rantai modul-$\Z$
			bebas $(C_n,\partial_n)_{n\in\Z}$ eksak dan
			$Z_n:=\Ker(\partial_n)$, maka barisan eksak pendek
			$0\to Z_{n+1}\to C_{n+1}\xrightarrow{\partial_n}Z_n\to0$
			terbelah untuk semua $n$.
		\end{hint}
		\item Buktikan bahwa setiap $P_n^*$ adalah modul-$G$ projektif yang
		dibangkitkan secara hingga.
		\begin{hint}
			Cukup tunjukkan $\Z[G]^*\simeq\Z[G]$.
		\end{hint}
		\item Misalkan $Q$ modul-$G$ projektif yang dibangkitkan secara hingga.
		Untuk setiap modul-$G$ $M$, buktikan bahwa morfisme kanonik
		$\nu:(Q\otimes M)_G\to(Q\otimes M)^G$ dari
		\eqref{eqn:nu-coinv-inv} merupakan isomorfisme.
		\begin{hint}
			Gunakan dekomposisi jumlah langsung untuk mereduksi ke kasus
			$Q=\Z[G]$, lalu verifikasikan langsung. Agar lebih sederhana,
			automorfisme
			\[ \left(\sum_{g\in G}a_gg\right)\otimes x
			\mapsto\sum_{g\in G}a_g(g\otimes g^{-1}x) \]
			mengubah aksi diagonal $G$ pada $\Z[G]\otimes M$ menjadi perkalian
			kiri.
		\end{hint}

		\item Definisikan $P_{-n}:=P^*_{n-1}$ dan sambungkan $P$ dengan $P^*$
		menjadi kompleks rantai
		\[ \mathcal{P}:=[\cdots\to P_1\to P_0
		\xrightarrow{\epsilon^*\epsilon}P_{-1}\to P_{-2}\to\cdots]. \]
		Buktikan bahwa untuk setiap modul-$G$ $M$ terdapat isomorfisme kanonik
		\[ \Hm^n\Hom_G(\mathcal{P},M)
		\simeq\text{kohomologi Tate }\TaHm^n(G,M),\qquad n\in\Z. \]

		\begin{hint}
			Untuk $n\geq1$ pernyataannya jelas. Untuk $n\leq-2$, gunakan sifat
			berikut: jika $Q$ adalah modul-$G$ projektif yang dibangkitkan secara
			hingga, maka $Q\otimes M\rightiso\Hom(Q^*,M)$. Dari sini deduksikan
			\[ (Q\otimes M)_G\xrightarrow[\nu]{\sim}(Q\otimes M)^G
			\rightiso\Hom(Q^*,M)^G=\Hom_G(Q^*,M). \]
			Untuk $n=0,-1$, identifikasikan
			$\Hom_G(P_{-1},M)\to\Hom_G(P_0,M)$ melalui isomorfisme di atas dengan
			\[ (P_0\otimes M)_G\xrightarrow{(\epsilon\otimes\identity)_G}M_G
			\xrightarrow{\nu}M^G\xrightarrow{\epsilon^*}\Hom_G(P_0,M). \]
		\end{hint}

		\item Gunakan hasil ini untuk membuktikan langsung barisan eksak panjang
		kohomologi Tate.
		\item Ambil resolusi bebas $\epsilon:P\to\Z$ dari
		Proposisi~\ref{prop:trivial-mod-nor-resolution}, pilih basis alaminya
		untuk setiap $P_n$, lalu deskripsikan secara eksplisit kompleks rantai
		$\mathcal{P}$ dari~(iv).
	\end{enumerate}

	Berdasarkan~(vi), hasil kali cawan dapat didefinisikan pada kohomologi
	Tate seperti dalam Lema~\ref{prop:cup-resolution}, dengan mendefinisikan
	$\Delta:\mathcal{P}\to\mathcal{P}\otimes\mathcal{P}$ secara sesuai.
	Untuk rumus eksplisit, lihat \cite[Bab IV, \S 7]{CF67}.
	\index{hasil kali cawan}
	\index{kohomologi Tate}
	\index[sym1]{HnGMTate}

	% segment-id: o014.li2.chapter6.exercises.ex018
	\item Misalkan $G$ grup hingga dan $H$ subgrupnya. Untuk modul-$G$ $M$,
	perluas melalui pergeseran dimensi (Korolari~\ref{prop:Tate-dim-shift})
	morfisme kanonik dalam \S\ref{sec:group-finite-index},
	$\mathrm{res}^n:\TaHm^n(G,M)\to\TaHm^n(H,M)$ dan
	$\mathrm{res}_n:\TaHm^{-n-1}(H,M)\to\TaHm^{-n-1}(G,M)$, dari kasus
	$n\geq1$ ke semua $n\in\Z$. Untuk sementara tuliskan yang pertama sebagai
	$\mathrm{res}_+$ dan yang kedua sebagai $\mathrm{res}_-$.
	\begin{enumerate}[(i)]
		\item Deskripsikan aksi $\mathrm{res}_+$ pada $\TaHm^0$ dan aksi
		$\mathrm{res}_-$ pada $\TaHm^{-1}$.
		\item Buktikan bahwa $\mathrm{res}_+$ pada $\TaHm^{-1}$ diinduksi oleh
		$\nu_{G|H}:M_G\to M_H$ (Definisi~\ref{def:cor-zero}), dan pada bagian
		$\TaHm^{<-1}$ sama dengan $\mathrm{cor}_n$
		(Definisi--Proposisi~\ref{def:cor}).
		\item Buktikan bahwa $\mathrm{res}_-$ pada $\TaHm^0$ diinduksi oleh
		$\nu^{G|H}:M^H\to M^G$, dan pada bagian $\TaHm^{>0}$ sama dengan
		$\mathrm{cor}^n$.
	\end{enumerate}

	% segment-id: o014.li2.chapter6.exercises.ex019
	\item Untuk grup $G$ dan kompleks modul-$G$ terbatas ke bawah
	$M=[\cdots\to M^p\to M^{p+1}\to\cdots]$, biarkan $p$ bervariasi dalam
	kompleks standar $C(G,M^p)$ untuk membentuk bikompleks. Tuliskan kompleks
	totalnya sebagai
	$C(G,M):=\tot((C^q(G,M^p))_{p,q})$. Buktikan isomorfisme kanonik
	\[ \Hm^n(C(G,M))\simeq
	\Hm^n(G,[\cdots\to M^p\to M^{p+1}\to\cdots]),\qquad n\in\Z, \]
	dengan ruas kanan didefinisikan seperti dalam \S\ref{sec:derived-primer}
	dan juga disebut hiperkohomologi grup.

	\begin{hint}
		Menurut definisi, ruas kanan diperoleh dengan mengambil resolusi injektif
		$M\to I$ (ingat: $I$ adalah kompleks terbatas ke bawah, setiap $I^q$
		modul-$G$ injektif, dan $M\to I$ kuasi-isomorfisme), lalu mengambil
		kohomologi dari kompleks $\Hom$
		$\Hom^\bullet_G(\Bbbk,I)$. Tunjukkan bahwa kohomologi yang sama dapat
		dihitung oleh $\Hom^\bullet_G(\mathsf{L},M)$.
	\end{hint}

	% segment-id: o014.li2.chapter6.exercises.ex020
	\item Buktikan hasil yang bersesuaian ketika $G$ grup profinit, dengan
	setiap $M^p$ disyaratkan modul-$G$ mulus dan kompleks standar kontinu
	$C(G,M^p)$ digunakan.

	\begin{hint}
		Definisikan funktor $C(G,\cdot)$ dari
		$\cate{C}^+(G\dcate{Mod}^\infty)$ ke
		$\cate{C}^+(\Bbbk\dcate{Mod})$. Tunjukkan bahwa funktor ini menginduksi
		funktor triangulasi pada tingkat $\cate{K}^+(\cdots)$ dan
		mempertahankan objek asiklik, sehingga menginduksi funktor triangulasi
		pada tingkat $\cate{D}^+(\cdots)$, tetap ditulis $C(G,\cdot)$. Melalui
		transformasi alami $(\cdot)^G\to C(G,\cdot)$ dan sifat universal
		funktor turunan kanan, diperoleh morfisme funktor triangulasi
		$\mathrm{R}(\cdot)^G\to C(G,\cdot)$. Untuk membuktikan bahwa ini
		isomorfisme, gunakan lema jalan keluar (Proposisi~\ref{prop:wayout})
		untuk mereduksi ke kasus Teorema~\ref{prop:profinite-group-standard}.
	\end{hint}

	% segment-id: o014.li2.chapter6.exercises.ex021
	\item Untuk grup profinit $G$ dan modul-$G$ hingga $M$, buktikan secara
	lengkap korespondensi bijektif yang disebutkan dalam
	\S\ref{sec:profinite-cohomology}, antara kelas ekuivalensi ekstensi grup
	profinit dan $\Hm^2(G,M)$.

	\begin{hint}
		Ikuti argumen dalam \S\ref{sec:grp-low-degree}. Dalam kasus grup profinit,
		pengamatan yang diperlukan adalah: (i) struktur grup yang didefinisikan
		oleh setiap $f\in\overline{Z}^2(G,M)$ pada ruang hasil kali $M\times G$
		merupakan grup profinit; (ii) untuk ekstensi grup profinit $E$, terdapat
		isomorfisme grup topologis $E/M\rightiso G$; (iii) dalam keadaan ini,
		Lema~\ref{prop:profinite-section} memberikan penampang kontinu
		$s:E\twoheadrightarrow G$ yang dapat disesuaikan agar ternormalisasi.
	\end{hint}

	% segment-id: o014.li2.chapter6.exercises.ex022
	\item Untuk medan $F$ sembarang, definisikan $\GL(n,F)$ sebagai grup
	matriks invertibel $n\times n$ di atas $F$, dan $\SL(n,F)$ sebagai subgrup
	yang ditentukan oleh $\det=1$. Perhatikan bahwa
	$G:=\Gal(\CC|\R)$ beraksi pada $\GL(n,\CC)$ dan $\SL(n,\CC)$, dengan
	titik-titik tetap masing-masing $\GL(n,\R)$ dan $\SL(n,\R)$.
	\begin{enumerate}[(i)]
		\item Ambil $n=2$. Misalkan $B:=\SL(2,\CC)$ beraksi melalui konjugasi
		pada ruang semua matriks $2\times2$, dan definisikan
		$A:=\Stab(\begin{smallmatrix}0&-1\\1&0\end{smallmatrix})$. Tunjukkan
		bahwa subgrup ini invarian di bawah aksi $G$ dan deskripsikan $A^G$.
		\item Tunjukkan bahwa $B/A$ (atau $B^G/A^G$) secara alami diidentifikasi
		dengan orbit $\left(\begin{smallmatrix}0&-1\\1&0\end{smallmatrix}\right)$
		di bawah aksi konjugasi $\SL(2,\CC)$ (atau $\SL(2,\R)$). Untuk $B/A$,
		jelaskan bagaimana aksi $G$ tercermin pada orbit tersebut.
		\item Tunjukkan $(B/A)^G\neq B^G/A^G$. Dengan kata lain, terdapat
		matriks real $2\times2$ yang terkonjugasi dengan
		$\left(\begin{smallmatrix}0&-1\\1&0\end{smallmatrix}\right)$ melalui
		$\SL(2,\CC)$ tetapi tidak melalui $\SL(2,\R)$.
		\item Tunjukkan bahwa fenomena di atas tidak terjadi jika $\SL$ diganti
		dengan $\GL$.
	\end{enumerate}

	% segment-id: o014.li2.chapter6.exercises.ex023
	\item Nyatakan dan buktikan kompatibilitas antara isomorfisme dalam
	Teorema~\ref{prop:Shapiro-noncommutative} dan homomorfisme penghubung dalam
	Teorema~\ref{prop:nonabelian-long}.

	% segment-id: o014.li2.chapter6.exercises.ex024
	\item (``Pemuntiran'' kohomologi grup) Gunakan notasi
	\S\ref{sec:nonabelian-coh}. Misalkan $A$ grup-$G$.
	\begin{enumerate}[(i)]
		\item Untuk $c\in Z^1(G,A)$, definisikan aksi grup baru
		$G\times A\to A$ dengan
		\[ {}^{c,g}a:=c(g)\cdot{}^ga\cdot c(g)^{-1},
		\qquad g\in G,\ a\in A. \]
		Tunjukkan bahwa ini memang memberikan grup-$G$, yang ditulis ${}^cA$.
		Buktikan pula bahwa jika $c\sim c'$, terdapat isomorfisme grup-$G$
		${}^{c'}A\simeq{}^cA$; deskripsikan secara eksplisit.
		\item Buktikan bahwa setiap $c\in Z^1(G,A)$ menginduksi bijeksi
		\begin{align*}
			\nu_c:Z^1(G,{}^cA)&\xrightarrow{1:1}Z^1(G,A),\\
			c'&\longmapsto c'c\quad\text{(perkalian titik demi titik)},
		\end{align*}
		yang menginduksi bijeksi
		${\bm\nu}_c:\Hm^1(G,{}^cA)\xrightarrow{1:1}\Hm^1(G,A)$ dengan titik
		pangkal dipetakan ke $[c]$. Berikan deskripsi yang disederhanakan ketika
		$A$ abelian.
		\item Misalkan terdapat barisan eksak pendek
		$1\to A\to B\to B/A\to1$ seperti dalam
		Teorema~\ref{prop:nonabelian-long}~(i). Untuk $c\in Z^1(G,A)$,
		buktikan diagram berikut komutatif:
		\[\begin{tikzcd}
			\Hm^1(G,{}^cA) \arrow[r] \arrow[d,"{{\bm\nu}_c}"'] &
			\Hm^1(G,{}^cB) \arrow[d,"{\bm{\nu}_c}"]\\
			\Hm^1(G,A) \arrow[r] & \Hm^1(G,B).
		\end{tikzcd}\]
		Gunakan ini untuk mendeskripsikan semua serat pemetaan
		$\Hm^1(G,A)\to\Hm^1(G,B)$.
		\item Nyatakan modifikasi yang diperlukan ketika $G$ grup profinit.
	\end{enumerate}
	\index{kohomologi grup!pemuntiran (twist)}
\end{Exercises}
